% Ad M. Brutum 1.9 --- headnote
% ad-brutum-01-09, early July 43 BC, Rome

\textit{Cicero to M. Brutus, from Rome, early July 43 BC
--- Perseus dateline \textit{Scr.~Romae.~in.~Quint.~...ut
videtur a.~711 (43)}, i.e. ``written at Rome at the
beginning of July (Quintilis), 43 BC, as it seems.'' The
\texttt{meta/works.yaml} entry carries the placeholder
date \texttt{-0043-02-26} at year precision, which is
plainly wrong: the dateline's \textit{in.~Quint.} places
the letter at the start of July, with the rest of the
1.9--1.14 cluster. The Perseus date is used in the
parallel sidecar (year-precision: \texttt{-0043-07}).}

\textit{This is a letter of condolence. Brutus's wife
Porcia --- Cato's daughter, and a figure of more than
ordinary moral weight in Brutus's circle --- has died.
Cicero turns the office of consolation around: when
\textit{he} lost Tullia, Brutus had been his consoler,
and stern with him when his grief was thought
unbecoming. Now Cicero returns the obligation, but with
a Roman insistence that Brutus, of all men, must be
seen to bear his loss with composure. The argument of
section 2 is the central note: ``my duty was to office
and to nature only; for you now it is to the people and
to the stage, as the saying goes, that you must do
service.'' \textit{Populo et scaenae} --- ``the people
and the stage'' --- the cause is publicly visible, the
liberator must perform the role. Section 3 closes
abruptly, with the political agenda surfacing through
the grief: ``We are awaiting you and your army.'' The
fuller dispatch on the state of the commonwealth is to
follow by ``our old friend'' --- the \textit{vetus
noster} who reappears as the courier in 1.12.}
